% Introduction — pivoted framing (the DSL kernel *evaluation* problem).
% Canonical: ../PIVOT_FRAMING.md. Pre-pivot gap-study drafts are in git history.
% TODO(revision): add references.bib entries for the \cite{kernelbench} key (KernelBench,
% Ouyang et al., Stanford ScalingIntelligence) and, if used, an LLM-kernel-generation
% reward-hacking example. \cite{kernelbench} is currently undefined (non-fatal build warning).

Deep learning systems are now widely deployed, and their efficiency depends heavily on a small number of GPU kernels that dominate execution time~\cite{acceleration_survey, optimization_survey}.
Frameworks have historically delegated these kernels to vendor libraries such as cuBLAS and cuDNN~\cite{cublas2024, chetlur2014cudnn}, but modern architectures increasingly require fused operators, specialized kernels, and model-specific computations that fall outside the fixed functionality those libraries expose~\cite{kernelfusion}.
Achieving high performance therefore now depends on the ability to generate efficient \emph{custom} kernels, and domain-specific languages (DSLs) such as Triton~\cite{openai2021triton} and TileLang~\cite{wang2025tilelang} have become the standard way to write them: they express kernels at the tile level while delegating memory movement, scheduling, and code generation to a compiler.
Their reach extends well beyond hand-written code---Triton is the lowering target for \texttt{torch.compile}~\cite{li2023torchcompile}, and DSLs are increasingly the output of LLM-based kernel generators~\cite{kernelbench}---so the efficiency of DSL-generated kernels is now a first-order concern.

The difficulty is that a developer, or an automated generator, who produces a \emph{functionally correct} DSL kernel has no reliable way to tell whether it is \emph{performant}.
A correct kernel may be competitive with the vendor library, or it may be orders of magnitude slower, and nothing in the source makes the difference visible.
When a DSL-generated kernel is correct yet much slower than the corresponding library routine, the developer must determine whether the loss arises from the algorithm, the schedule, or the compiler---and current toolchains offer little diagnostic support for that distinction~\cite{gpuenergy}.
This makes performance deficits hard to predict and hard to correct, which is a software-engineering problem as much as a compiler one: a high-level kernel language is useful in practice only if its performance behavior is something developers can evaluate and reason about.

The benchmarks practitioners actually rely on do not close this gap.
The prevailing benchmarks for DSL and LLM-generated kernels---KernelBench~\cite{kernelbench} and TritonBench~\cite{li2025tritonbench}---gate on \emph{correctness}, reporting performance as an outcome rather than enforcing it as an acceptance criterion.
As a result they admit performance-poor kernels: in our measurements a naive but idiomatic TileLang kernel passes the KernelBench-style correctness gate while running more than $300\times$ slower than PyTorch, and the gate's one reward-hacking guard never fires because it only flags suspiciously \emph{fast} kernels.
These benchmarks also cover a narrow slice of the space---often a single DSL, data type, shape, and GPU per task---so a kernel that passes tells a developer little about whether it is well written.
And the obvious remedy, a benchmark comprehensive enough to span every operator, data type, shape, and architecture, is combinatorially infeasible to build and maintain.

This paper studies that evaluation problem directly.
We ask three questions: (RQ1) whether existing benchmarks can distinguish well-written DSL kernels from performance-poor ones, and where they fall short; (RQ2) for kernels that pass such benchmarks, how large and structured the hidden performance gap to vendor libraries actually is, and what causes it; and (RQ3) how to guide DSL kernel development when a comprehensive benchmark is infeasible.
We evaluate Triton and TileLang against PyTorch with cuBLAS/cuDNN on 22 kernels spanning five operator classes---GEMM, attention, convolution, normalization, and element-wise/reduction---with hardware-counter profiling via Nsight Compute.
Primary measurements span two datacenter architectures---the NVIDIA A100-SXM4-40GB (Ampere) and the NVIDIA GH200 (Grace Hopper).

Our findings are threefold.
First, the evaluation gap is real and large: naive-but-correct kernels pass the gate while running hundreds of times slower---the shipped TileLang LayerNorm is $323\times$ slower under locked clocks---and the dominant benchmarks neither gate on performance nor cover enough of the space to catch it.
Second, the hidden gap is highly structured: most is an \emph{authoring} artifact that a single reduction-primitive change (\texttt{T.serial}$\rightarrow$\texttt{T.reduce} with native-dtype I/O) repairs to library parity or better, while a residual convolution gap is structural (cuDNN's implicit-GEMM maturity; absent Winograd contributes only $\approx$2--3\%); we separate causes by \emph{where the fix belongs}---user-space authoring, code generation, and library maturity.
Third, absent a comprehensive benchmark, lightweight heuristics suffice: a library-comparability check anchored by a baseline-independent roofline fraction cleanly separates well-written kernels from poor ones, and a small set of recurring optimization patterns guides the fixes.

\smallskip\noindent\textbf{Contributions.}
\begin{itemize}
    \item \textbf{The evaluation gap (RQ1).} We show that prevailing DSL and LLM-generated-kernel benchmarks admit performance-poor kernels---demonstrated end-to-end against the actual evaluator---and survey their coverage and reward-hacking gaps.
    \item \textbf{The hidden gap and its causes (RQ2).} We characterize the performance gap these benchmarks miss across 22 kernels and five operator classes on multiple GPU architectures, with hardware-counter root causes that separate authoring, code-generation, and library-maturity factors.
    \item \textbf{Guidance without a benchmark (RQ3).} We propose and validate lightweight evaluation heuristics---a library-comparability screen plus a baseline-independent roofline anchor---and a set of recurring optimization patterns that guide DSL kernel development and recover most of the gap.
    \item \textbf{Artifact.} We release the kernel suite, profiling infrastructure, and analysis to support reproducibility and follow-on work.
\end{itemize}
